% --- Archon physics formalization source begin ---
% archon:physics
% archon:covers PhyXMiniProblems/problem_phyx_mini_0574.lean
% archon:source-report reports/phyx_mini/problem_phyx_mini_0574.source.json

\chapter{Physics problem phyx\_mini\_0574}
\label{ch:PhyXMiniProblems_problem_phyx_mini_0574}

\paragraph{Problem source.}
\#\# Physical scenario

Figure shows part of the decay scheme of \$\textasciicircum{}\{237\}\textbackslash{}text\{Np\}\$ on a plot of mass number \$A\$ versus proton number \$Z\$; five lines that represent either alpha decay or beta-minus decay connect dots that represent isotopes.

\#\# Question

What is the isotope at the end of the five decays (as marked with a question mark in figure)?

\#\# Answer choices

- A: A: \{\}\textasciicircum{}\{225\}Ra

- B: B: \{\}\textasciicircum{}\{229\}Ac

- C: C: \{\}\textasciicircum{}\{225\}Ac

- D: D: \{\}\textasciicircum{}\{223\}Th

\#\# Figure caption (auxiliary; use the image as primary evidence)

The image is a section of a graph with two axes labeled "A" (vertical axis) and "Z" (horizontal axis). It appears to represent nuclear properties or isotopes.

- The graph has a series of connected black dots linked by a green line, forming a zigzag path.
- The dots are plotted on a grid. 
- The path starts at a lower-left point and moves to the upper-right.
- The path ends at a point labeled with the isotope \textbackslash{}( \textasciicircum{}\{237\}\textbackslash{}text\{Np\} \textbackslash{}) on the top right.
- There is a question mark next to the second point on the path, which suggests an unknown or unspecified isotope at that position.
- The axes have a break symbol near the origin, indicating a discontinuity or a skipped range of values.

The graph likely depicts a series of nuclear reactions or decay paths.

\#\# Dataset metadata

Nuclear Physics; Physical Model Grounding Reasoning, Multi-Formula Reasoning

\paragraph{Recorded answer/context.}
C: C: \{\}\textasciicircum{}\{225\}Ac

\paragraph{Figure/image path.}
/root/proposal\_for\_physic/science-mango/phyx\_mini\_run/phyx\_data/test\_image/574.png

\paragraph{Formalization target.}
create a compiling Lean file with sorry bodies at `PhyXMiniProblems/problem\_phyx\_mini\_0574.lean`. The Lean declarations must preserve the physical quantities, dimensions or dimensional roles, figure labels, governing-law hypotheses, and final relation expressed by this problem.
Use Mathlib/Physlib names found through LeanExplore where available. If a physics API is missing, introduce faithful local abstractions rather than scalar placeholder aliases.

\begin{theorem}[Physics formalization target]
\label{thm:physics:phyx_mini_0574:target}
\lean{PhyXMiniProblems.ProblemPhyXMini0574.questionMark_isotope_is_choice_C}
\uses{def:physics:phyx-mini-0574:phyxminiproblems-problemphyxmini0574-answerisotope, def:physics:phyx-mini-0574:phyxminiproblems-problemphyxmini0574-fivestepdecayscheme, def:physics:phyx-mini-0574:phyxminiproblems-problemphyxmini0574-matchessupplieddecayfigure, def:physics:phyx-mini-0574:phyxminiproblems-problemphyxmini0574-satisfiesdecaynumberlaws}
The assigned autoformalize agent should translate this physics problem into Lean declarations in the covered file, with theorem and lemma proof bodies written as `by sorry`.
\end{theorem}
\begin{proof}
This is an autoformalization task, not a proof task. Produce statements that can later be proved by the physics prover without weakening the physical model.
\end{proof}

% --- Archon declaration topology begin ---
\section{Declaration topology}
\begin{definition}[Chemical Element]
  \label{def:physics:phyx-mini-0574:phyxminiproblems-problemphyxmini0574-chemicalelement}
  \lean{PhyXMiniProblems.ProblemPhyXMini0574.ChemicalElement}
  Chemical elements occurring in the displayed decay-chain segment or choices
\end{definition}
\begin{proof}
  This is the declared model definition of Chemical Element.
\end{proof}
\begin{definition}[atomic Number]
  \label{def:physics:phyx-mini-0574:phyxminiproblems-problemphyxmini0574-chemicalelement-atomicnumber}
  \lean{PhyXMiniProblems.ProblemPhyXMini0574.ChemicalElement.atomicNumber}
  \uses{def:physics:phyx-mini-0574:phyxminiproblems-problemphyxmini0574-chemicalelement}
  Proton number Z of each chemical element relevant to this problem
\end{definition}
\begin{proof}
  This is the declared model definition of atomic Number.
\end{proof}
\begin{definition}[Nuclear Isotope]
  \label{def:physics:phyx-mini-0574:phyxminiproblems-problemphyxmini0574-nuclearisotope}
  \lean{PhyXMiniProblems.ProblemPhyXMini0574.NuclearIsotope}
  \uses{def:physics:phyx-mini-0574:phyxminiproblems-problemphyxmini0574-chemicalelement}
  A nuclide specified by mass number A and chemical element. Its proton number is the atomic number of the element; the inequality excludes nuclides with more protons than nucleons
\end{definition}
\begin{proof}
  This is the declared model definition of Nuclear Isotope.
\end{proof}
\begin{definition}[proton Number]
  \label{def:physics:phyx-mini-0574:phyxminiproblems-problemphyxmini0574-nuclearisotope-protonnumber}
  \lean{PhyXMiniProblems.ProblemPhyXMini0574.NuclearIsotope.protonNumber}
  \uses{def:physics:phyx-mini-0574:phyxminiproblems-problemphyxmini0574-nuclearisotope}
  Proton-number coordinate Z of a nuclide on the supplied plot
\end{definition}
\begin{proof}
  This is the declared model definition of proton Number.
\end{proof}
\begin{definition}[neptunium237]
  \label{def:physics:phyx-mini-0574:phyxminiproblems-problemphyxmini0574-neptunium237}
  \lean{PhyXMiniProblems.ProblemPhyXMini0574.neptunium237}
  \uses{def:physics:phyx-mini-0574:phyxminiproblems-problemphyxmini0574-nuclearisotope}
  The labelled source nuclide, neptunium-237
\end{definition}
\begin{proof}
  This is the declared model definition of neptunium237.
\end{proof}
\begin{definition}[radium225]
  \label{def:physics:phyx-mini-0574:phyxminiproblems-problemphyxmini0574-radium225}
  \lean{PhyXMiniProblems.ProblemPhyXMini0574.radium225}
  \uses{def:physics:phyx-mini-0574:phyxminiproblems-problemphyxmini0574-nuclearisotope}
  Radium-225, displayed as answer choice A
\end{definition}
\begin{proof}
  This is the declared model definition of radium225.
\end{proof}
\begin{definition}[actinium229]
  \label{def:physics:phyx-mini-0574:phyxminiproblems-problemphyxmini0574-actinium229}
  \lean{PhyXMiniProblems.ProblemPhyXMini0574.actinium229}
  \uses{def:physics:phyx-mini-0574:phyxminiproblems-problemphyxmini0574-nuclearisotope}
  Actinium-229, displayed as answer choice B
\end{definition}
\begin{proof}
  This is the declared model definition of actinium229.
\end{proof}
\begin{definition}[actinium225]
  \label{def:physics:phyx-mini-0574:phyxminiproblems-problemphyxmini0574-actinium225}
  \lean{PhyXMiniProblems.ProblemPhyXMini0574.actinium225}
  \uses{def:physics:phyx-mini-0574:phyxminiproblems-problemphyxmini0574-nuclearisotope}
  Actinium-225, displayed as answer choice C
\end{definition}
\begin{proof}
  This is the declared model definition of actinium225.
\end{proof}
\begin{definition}[thorium223]
  \label{def:physics:phyx-mini-0574:phyxminiproblems-problemphyxmini0574-thorium223}
  \lean{PhyXMiniProblems.ProblemPhyXMini0574.thorium223}
  \uses{def:physics:phyx-mini-0574:phyxminiproblems-problemphyxmini0574-nuclearisotope}
  Thorium-223, displayed as answer choice D
\end{definition}
\begin{proof}
  This is the declared model definition of thorium223.
\end{proof}
\begin{definition}[Answer Choice]
  \label{def:physics:phyx-mini-0574:phyxminiproblems-problemphyxmini0574-answerchoice}
  \lean{PhyXMiniProblems.ProblemPhyXMini0574.AnswerChoice}
  Labels of the four displayed multiple-choice answers
\end{definition}
\begin{proof}
  This is the declared model definition of Answer Choice.
\end{proof}
\begin{definition}[answer Isotope]
  \label{def:physics:phyx-mini-0574:phyxminiproblems-problemphyxmini0574-answerisotope}
  \lean{PhyXMiniProblems.ProblemPhyXMini0574.answerIsotope}
  \uses{def:physics:phyx-mini-0574:phyxminiproblems-problemphyxmini0574-nuclearisotope, def:physics:phyx-mini-0574:phyxminiproblems-problemphyxmini0574-radium225, def:physics:phyx-mini-0574:phyxminiproblems-problemphyxmini0574-actinium229, def:physics:phyx-mini-0574:phyxminiproblems-problemphyxmini0574-actinium225, def:physics:phyx-mini-0574:phyxminiproblems-problemphyxmini0574-thorium223, def:physics:phyx-mini-0574:phyxminiproblems-problemphyxmini0574-answerchoice}
  Isotope printed beside each answer-choice label
\end{definition}
\begin{proof}
  This is the declared model definition of answer Isotope.
\end{proof}
\begin{definition}[Plot Axis]
  \label{def:physics:phyx-mini-0574:phyxminiproblems-problemphyxmini0574-plotaxis}
  \lean{PhyXMiniProblems.ProblemPhyXMini0574.PlotAxis}
  Geometric axes of the supplied decay-scheme plot
\end{definition}
\begin{proof}
  This is the declared model definition of Plot Axis.
\end{proof}
\begin{definition}[Nuclear Plot Quantity]
  \label{def:physics:phyx-mini-0574:phyxminiproblems-problemphyxmini0574-nuclearplotquantity}
  \lean{PhyXMiniProblems.ProblemPhyXMini0574.NuclearPlotQuantity}
  Nuclear count plotted on an axis; both quantities are dimensionless
\end{definition}
\begin{proof}
  This is the declared model definition of Nuclear Plot Quantity.
\end{proof}
\begin{definition}[physical Dimension]
  \label{def:physics:phyx-mini-0574:phyxminiproblems-problemphyxmini0574-nuclearplotquantity-physicaldimension}
  \lean{PhyXMiniProblems.ProblemPhyXMini0574.NuclearPlotQuantity.physicalDimension}
  \uses{def:physics:phyx-mini-0574:phyxminiproblems-problemphyxmini0574-nuclearplotquantity}
  Physical dimension of either nuclear count plotted in the figure. Physlib's multiplicative identity dimension records that A and Z are dimensionless; their integer readouts remain natural numbers
\end{definition}
\begin{proof}
  This is the declared model definition of physical Dimension.
\end{proof}
\begin{definition}[Decay Path Point]
  \label{def:physics:phyx-mini-0574:phyxminiproblems-problemphyxmini0574-decaypathpoint}
  \lean{PhyXMiniProblems.ProblemPhyXMini0574.DecayPathPoint}
  The six black dots, ordered from neptunium-237 toward the question mark
\end{definition}
\begin{proof}
  This is the declared model definition of Decay Path Point.
\end{proof}
\begin{definition}[Decay Step]
  \label{def:physics:phyx-mini-0574:phyxminiproblems-problemphyxmini0574-decaystep}
  \lean{PhyXMiniProblems.ProblemPhyXMini0574.DecayStep}
  The five green line segments, ordered in the physical decay direction
\end{definition}
\begin{proof}
  This is the declared model definition of Decay Step.
\end{proof}
\begin{definition}[start Point]
  \label{def:physics:phyx-mini-0574:phyxminiproblems-problemphyxmini0574-decaystep-startpoint}
  \lean{PhyXMiniProblems.ProblemPhyXMini0574.DecayStep.startPoint}
  \uses{def:physics:phyx-mini-0574:phyxminiproblems-problemphyxmini0574-decaypathpoint, def:physics:phyx-mini-0574:phyxminiproblems-problemphyxmini0574-decaystep}
  Initial plotted point of a green decay segment
\end{definition}
\begin{proof}
  This is the declared model definition of start Point.
\end{proof}
\begin{definition}[end Point]
  \label{def:physics:phyx-mini-0574:phyxminiproblems-problemphyxmini0574-decaystep-endpoint}
  \lean{PhyXMiniProblems.ProblemPhyXMini0574.DecayStep.endPoint}
  \uses{def:physics:phyx-mini-0574:phyxminiproblems-problemphyxmini0574-decaypathpoint, def:physics:phyx-mini-0574:phyxminiproblems-problemphyxmini0574-decaystep}
  Final plotted point of a green decay segment
\end{definition}
\begin{proof}
  This is the declared model definition of end Point.
\end{proof}
\begin{definition}[Decay Mode]
  \label{def:physics:phyx-mini-0574:phyxminiproblems-problemphyxmini0574-decaymode}
  \lean{PhyXMiniProblems.ProblemPhyXMini0574.DecayMode}
  The two kinds of radioactive decay stipulated for the plotted segments
\end{definition}
\begin{proof}
  This is the declared model definition of Decay Mode.
\end{proof}
\begin{definition}[Decay Scheme Figure]
  \label{def:physics:phyx-mini-0574:phyxminiproblems-problemphyxmini0574-decayschemefigure}
  \lean{PhyXMiniProblems.ProblemPhyXMini0574.DecaySchemeFigure}
  \uses{def:physics:phyx-mini-0574:phyxminiproblems-problemphyxmini0574-plotaxis, def:physics:phyx-mini-0574:phyxminiproblems-problemphyxmini0574-nuclearplotquantity, def:physics:phyx-mini-0574:phyxminiproblems-problemphyxmini0574-decaypathpoint, def:physics:phyx-mini-0574:phyxminiproblems-problemphyxmini0574-decaystep}
  Raster-level content of the supplied A-versus-Z diagram
\end{definition}
\begin{proof}
  This is the declared model definition of Decay Scheme Figure.
\end{proof}
\begin{definition}[Five Step Decay Scheme]
  \label{def:physics:phyx-mini-0574:phyxminiproblems-problemphyxmini0574-fivestepdecayscheme}
  \lean{PhyXMiniProblems.ProblemPhyXMini0574.FiveStepDecayScheme}
  \uses{def:physics:phyx-mini-0574:phyxminiproblems-problemphyxmini0574-nuclearisotope, def:physics:phyx-mini-0574:phyxminiproblems-problemphyxmini0574-decaypathpoint, def:physics:phyx-mini-0574:phyxminiproblems-problemphyxmini0574-decaystep, def:physics:phyx-mini-0574:phyxminiproblems-problemphyxmini0574-decaymode, def:physics:phyx-mini-0574:phyxminiproblems-problemphyxmini0574-decayschemefigure}
  Independent quantities in the five-step scheme. In particular, the isotope at the question mark is an unconstrained field until figure evidence and the general decay laws are supplied
\end{definition}
\begin{proof}
  This is the declared model definition of Five Step Decay Scheme.
\end{proof}
\begin{definition}[Matches Supplied Decay Figure]
  \label{def:physics:phyx-mini-0574:phyxminiproblems-problemphyxmini0574-matchessupplieddecayfigure}
  \lean{PhyXMiniProblems.ProblemPhyXMini0574.MatchesSuppliedDecayFigure}
  \uses{def:physics:phyx-mini-0574:phyxminiproblems-problemphyxmini0574-neptunium237, def:physics:phyx-mini-0574:phyxminiproblems-problemphyxmini0574-fivestepdecayscheme}
  Facts read from the primary image. The sloping down-left displacements in the decay direction are alpha steps, while the same-A, one-column-right displacements are beta-minus steps. No endpoint isotope is asserted here
\end{definition}
\begin{proof}
  This is the declared model definition of Matches Supplied Decay Figure.
\end{proof}
\begin{definition}[Is Decay Product]
  \label{def:physics:phyx-mini-0574:phyxminiproblems-problemphyxmini0574-isdecayproduct}
  \lean{PhyXMiniProblems.ProblemPhyXMini0574.IsDecayProduct}
  \uses{def:physics:phyx-mini-0574:phyxminiproblems-problemphyxmini0574-nuclearisotope, def:physics:phyx-mini-0574:phyxminiproblems-problemphyxmini0574-decaymode}
  Conservation-number effect of one decay. Alpha decay removes two protons and four nucleons from the parent; beta-minus decay preserves A and converts a neutron into a proton, increasing the daughter's Z by one
\end{definition}
\begin{proof}
  This is the declared model definition of Is Decay Product.
\end{proof}
\begin{definition}[Satisfies Decay Number Laws]
  \label{def:physics:phyx-mini-0574:phyxminiproblems-problemphyxmini0574-satisfiesdecaynumberlaws}
  \lean{PhyXMiniProblems.ProblemPhyXMini0574.SatisfiesDecayNumberLaws}
  \uses{def:physics:phyx-mini-0574:phyxminiproblems-problemphyxmini0574-fivestepdecayscheme, def:physics:phyx-mini-0574:phyxminiproblems-problemphyxmini0574-isdecayproduct}
  Every pictured segment obeys the number law for its classified mode
\end{definition}
\begin{proof}
  This is the declared model definition of Satisfies Decay Number Laws.
\end{proof}
\begin{theorem}[question Mark coordinates]
  \label{thm:physics:phyx-mini-0574:phyxminiproblems-problemphyxmini0574-questionmark-coordinates}
  \lean{PhyXMiniProblems.ProblemPhyXMini0574.questionMark_coordinates}
  \uses{def:physics:phyx-mini-0574:phyxminiproblems-problemphyxmini0574-fivestepdecayscheme, def:physics:phyx-mini-0574:phyxminiproblems-problemphyxmini0574-matchessupplieddecayfigure, def:physics:phyx-mini-0574:phyxminiproblems-problemphyxmini0574-satisfiesdecaynumberlaws}
  The five decay-number updates determine the two endpoint coordinates: mass number A = 225 and proton number Z = 89
\end{theorem}
\begin{proof}
  The conclusion follows from the governing relations and figure readouts named in the statement of question Mark coordinates.
\end{proof}
% --- Archon declaration topology end ---
% --- Archon physics formalization source end ---
